\section{Đặc tả dự án}

\subsection{Bối cảnh dự án}

Phương tiện cá nhân là công cụ hữu dụng bậc nhất đáp ứng nhu cầu đi lại đặc thù đối với một người. Dưới góc nhìn của một đơn vị giáo dục như Đại học Bách khoa, tổ chức các khu vực tập kết để gửi xe cho các cán bộ, giảng viên và sinh viên của nhà trường là điều tiên quyết nhằm tạo ra sự thuận lợi trong việc đi học và đi dạy. Tuy nhiên, đối mặt với quy mô và sự tăng trưởng về số lượng người tham gia học và dạy, phương thức quản lý thủ công hoặc sử dụng các công cụ rời rạc ngày càng bộc lộ nhiều hạn chế. Điều này dẫn đến những khó khăn trong việc theo dõi, phân luồng xe và thu phí gửi người ra vào khuôn viên trường. Vì vậy, trường Đại học Bách khoa mong muốn xây dựng một hệ thống phần mềm hiện đại, có khả năng mở rộng, nhằm quản lý và vận hành chương trình hiệu quả và bền vững, đồng thời đáp ứng nhu cầu thực tiễn trong môi trường giáo dục đại học.

Quy trình quản lý hiện nay tồn tại nhiều hạn chế, ảnh hưởng đến cả giảng viên, sinh viên và nhân viên bãi xe: sinh viên gặp khó khăn trong việc tìm kiếm chỗ đỗ xe; giảng viên chật vật lấy xe ra giữa các hàng xe đỗ thiếu quy chuẩn cũng như bị chiếm dụng chỗ đỗ bởi các cá nhân khác; nhân sự bãi xe gặp khó khăn trong việc sắp xếp bãi xe nhằm đáp ứng só lượng và chất lượng của chỗ đỗ.

Hệ thống phần mềm được kỳ vọng sẽ trở thành nền tảng phân luồng xe một cách thuận tiện, phù hợp đáp ứng nhu cầu về chỗ đỗ riêng biệt dành cho giảng viên và sinh viên. Hệ thống cũng sẽ tối ưu hóa quy trình quản lý và vận hành, mang lại trải nghiệm ra - vào bãi xe và thanh toán chi phí gửi xe thoải mái và thuận tiện.